% mage-auto-heading: lifecycle
% mage-auto-heading: ownership-and-lease
% Authoritative LaTeX manuscript; Markdown is a generated compatibility projection.

Behavioral models represent state, transition, and time. Ownership models represent who may act on or hold a resource. The two are distinct but often joined: behavior constrains possible transitions; ownership constrains who may perform them.

\section{Lifecycle}\label{heading:lifecycle}

\textbf{Engineering question.} What states may this object occupy, and which transitions are legal?

\textbf{Representation.} A minimal job lifecycle moves an item from \texttt{FREE} to \texttt{LEASED} on claim, from \texttt{LEASED} to \texttt{DONE} on completion, and from \texttt{LEASED} back to \texttt{FREE} on expiry or release. \texttt{DONE} is terminal.

\textbf{Property.} The invariants sort into safety and liveness:

{\footnotesize
\begin{longtable}[]{@{}
  >{\raggedright\arraybackslash}p{(\linewidth - 2\tabcolsep) * \real{0.5000}}
  >{\raggedright\arraybackslash}p{(\linewidth - 2\tabcolsep) * \real{0.5000}}@{}}
\caption[Lifecycle invariants]{Representative lifecycle invariants, by kind.}\tabularnewline
\toprule\noalign{}
\begin{minipage}[b]{\linewidth}\raggedright
Invariant
\end{minipage} & \begin{minipage}[b]{\linewidth}\raggedright
Kind
\end{minipage} \\
\midrule\noalign{}
\endfirsthead
\toprule\noalign{}
\begin{minipage}[b]{\linewidth}\raggedright
Invariant
\end{minipage} & \begin{minipage}[b]{\linewidth}\raggedright
Kind
\end{minipage} \\
\midrule\noalign{}
\endhead
\bottomrule\noalign{}
\endlastfoot
Only declared transitions occur & Safety \\
Terminal states do not re-enter processing & Safety \\
A claimed job eventually completes or returns to an available state, under stated assumptions & Liveness \\
\end{longtable}
}

\textbf{Enforcement and correspondence.} Transition operations or checks can mediate actual state changes against the declared transition relation. Any implemented transition absent from that relation becomes a finding. Where two lifecycles interact, represent their cross-machine transitions explicitly rather than leaving those relationships implicit in code.

\section{Ownership and Lease}\label{heading:ownership-and-lease}

Concurrency adds ownership and lease state to the lifecycle. The resulting model must represent not only which transitions are legal, but which actor may perform them now.

\textbf{Engineering question.} Who may act on this work item now, and what stays true if workers overlap, fail, or retry?

\autoref{fig-g3-behavior-ownership} lays the ownership overlay over the state machine.

\begin{figure}[htbp]
\centering
\includegraphics[width=\linewidth]{book/latex/assets/appendix-g-3-behavior-ownership-d07671b3d7.pdf}
\caption{\emph{Behavior plus ownership.} The lifecycle FREE → LEASED → DONE is combined with a lease that records the current owner and expiry. The ownership invariant requires exactly one valid owner while LEASED and no active owner while FREE.}
\label{fig-g3-behavior-ownership}
\end{figure}

\textbf{Property.} Representative invariants:

\begin{itemize}
\tightlist
\item
  \textbf{A leased item has a valid owner}, and a free item has none.
\item
  \textbf{Ownership changes through the declared acquisition and release mechanism}, not by ad hoc mutation.
\item
  \textbf{Terminal work cannot be reclaimed.}
\item
  \textbf{Where ordering is required, observed acquisition order respects the declared order.}
\end{itemize}

\textbf{Enforcement and correspondence.} Some ownership facts are runtime facts and must be observed. The ownership protocol itself is authored intent. Atomic operations, leases, compare-and-swap, transition primitives, and admission checks can make selected invariants enforceable.

Concurrency limits and ownership answer different questions. A semaphore or mediator constrains how many actors may execute; single-writer ownership identifies which actor may mutate. They require distinct representations because they govern different properties.
